\section{Computations}\label{app:data}

\subsection{The sheared curves}\label{ss:methods}
Apart from the base case of Remark~\ref{rem:chirality}, a verified computation described at the end of this
subsection, the proofs in Sections~\ref{sec:sheared} and~\ref{sec:twisted} do not use computation.  The following
computations corroborate them.
\begin{enumerate}
\item[(i)] \emph{Lemma~\ref{lem:straight}.}  Each displayed identity in Steps 1--7 of its proof holds
 identically in the indeterminates $v$ and $\zeta$, as verified by symbolic computation.  Without the
 assumption $E=0$ one has $\mathbf{Vec}\,\rho(w_0)\cdot\mathbf{Vec}\,\rho(a)=\frac12\zeta\sin2v\,E^2$ and
 $\operatorname{tr}\bigl(\rho(w_0)\rho(a)\bigr)=-2\zeta\sin2v\,E^2$.  In the Fricke coordinates
 $x=\operatorname{tr}\rho(A)$, $y=\operatorname{tr}\rho(B)$, $z=\operatorname{tr}\rho(AB)$, the condition
 $\operatorname{tr}\rho(b)=0$ reads $x=yz$; then $E=y^2+z^2-3$,
 \[
  \operatorname{tr}\rho(a)=-yE,\qquad \operatorname{tr}\rho(w_0)=-zE,\qquad \operatorname{tr}\rho(w_0a)=yzE^2,
 \]
 and $\operatorname{tr}\rho(w)-2$ lies in the ideal generated by $x-yz$ and $E$.  For the word
 $(ba)^{-2}c\,(ba)^2a$ the same reduction leaves $-y^2-1$, so that word is not trivial on $C_0$.  At three
 points of $C_0$ with algebraic coordinates, exact quaternion arithmetic gives $\rho(w)=1$, and a
 numerical trace of $\Pi(C_0)$ at $50$ significant digits satisfies $|\theta-6\gamma+\pi|\le1.4\cdot10^{-47}$.
\item[(ii)] \emph{Combinatorics.}  For every odd $q\le401$ we verified in exact rational arithmetic: the
 crossings of the sheared polygons of Corollary~\ref{cor:polygon} with the lifts of $L_0$, for both
 resolutions; the value of $N_q$; the order of the crossing abscissae in $(\frac\pi6,\frac{5\pi}6)$ and
 the absence of a crossing at $\gamma=\frac\pi2$; the bookkeeping of the graded homology in
 Theorem~\ref{thm:pillowP}; the count of Remark~\ref{rem:FKPcount}; and the number of generators on the
 remnant circle.
\item[(iii)] \emph{Groups.}  The identity of Lemma~\ref{lem:relator} holds as an identity of reduced free-group
 words for $n\in\{5,8,\dots,29\}$ and $m\le15$; with $t_e^{+m}$ in place of $t_e^{-m}$ the relator of~\eqref{eq:CW} is
 never obtained.  The boundary relation $ba=cd$ holds for the words \eqref{eq:words-n}, and for all $4\le n\le26$ prime
 to $3$ and $m\le6$ the Alexander polynomial obtained by Fox calculus from the presentation of
 Lemma~\ref{lem:group} with $h=-m$ equals that of the closure of $(\sigma_1\sigma_2)^n\sigma_1^{2m}$
 computed by the Burau representation.  The closed form of the Alexander polynomials in Proposition~\ref{prop:twisted-alex} was checked as an identity of rational functions and exactly for $1092$
 knots, and the determinants for $n\le50$ and $m\le20$.
\item[(iv)] \emph{The shear.}  Applying $t_e^h$ to the boundary values of points of the traced variety
 $W_\varepsilon$ reproduces $S_h$ to within $10^{-12}$ for $h=-6,\dots,2$, whereas the shear of the
 opposite sign is off by about $\pi$.
\item[(v)] \emph{Numerical Floer complexes.}  The variety $W_\varepsilon$ is traced with the continuation
 code of \cite{WuebbenV}.  Independent code applies $S_{-m}$, locates the generators on the lifts of
 $L_0$, computes their degrees from \eqref{eq:grading15}, and enumerates the bigons whose lifts are embedded
 discs, as Jordan domains in $\R^2$.  Maslov indices are computed on the unsheared curve against the line field spanned by
 $(1,6-q)$, which by Lemma~\ref{lem:shear}(f) gives the index of the sheared curve against $\ell_1$; an
 index obtained by unwrapping tangent angles along the sheared polygonal curve is unreliable near the
 junctions when $\varepsilon$ is small or $q$ is large.  The results for the pretzel knots, summarized in
 Remark~\ref{rem:numericsP}, are as follows.
 \begin{itemize}
 \item \emph{Joined sectors.} For ten perturbations with $|5\varepsilon_A|$ between $0.002$ and $0.3$ and
 $|\varepsilon_B|$ between $0.0003$ and $0.05$, two earring parameters, and $q\in\{7,9,\dots,25,29,33\}$, the
 homology agrees with Theorem~\ref{thm:pillowP} in all $240$ cases. Every bigon found runs from a down-crossing to
 an up-crossing on one lift of $L_0$, the number of bigons is twice the rank of the differential, and the degrees
 agree with Proposition~\ref{prop:degrees}. The generator count $q+2+4N_q$ is attained in $178$ cases; in the
 other $62$, at the larger perturbations, canceling pairs are missing, as in Remark~\ref{rem:threshold}.
 \item \emph{Split sectors.} For five perturbations with $|\varepsilon_B|>\sqrt3|\varepsilon_A|$ and
 $|\varepsilon|\le0.031$, and $q=7,9,\dots,21$, all $40$ cases agree: the arc is $\varphi$-monotone with $q-2$
 generators, and the remnant circle has homology of rank $1$ in each degree. The circle carries $4+4N_q$
 generators in $37$ cases; in the other three ($q=17$, at the three larger perturbations) a canceling pair is
 missing. A lift of $L_0$ meets the lift of the circle up to seven times ($q=21$).
 \item \emph{Resolution type.} At $|\varepsilon|=0.03$, in $24$ equally spaced directions of which four lie on
 the excluded lines, the observed type agrees with Proposition~\ref{prop:restype} in all $20$ remaining
 directions; the split ones are those at $75^\circ$, $90^\circ$, $105^\circ$, $255^\circ$, $270^\circ$ and
 $285^\circ$. The sample points of~\cite[Rem.~5.8]{WuebbenV} for $n=5$, in both kinds of sector, also agree with
 Proposition~\ref{prop:restype}.
 \end{itemize}
 For the twisted torus
 knots with $n=7,8,13$ and $m=1,2,3$, at two perturbations each, the graded homology agrees with
 Theorem~\ref{thm:twisted} in all $18$ runs.
\end{enumerate}

\emph{The base case of Section~\ref{sec:sheared}.}  For Remark~\ref{rem:chirality}, the complements of $K^{(0)}\cup e$ (from a
planar diagram of~\cite[Fig.~18]{HHK2}) and of the closure of $(\sigma_1\sigma_2)^5$ with an unknotted circle around two
strands are compared through their canonical cell decompositions, computed with interval arithmetic by SnapPy within
Sage~\cite{SnapPy,HIKMOT,DHL}.  The comparison lists all isometries between the two complements: there are two, and both
preserve the orientation and the meridians.  The surgery sign convention is checked by a Rolfsen twist: $-1$ surgery
on the circle gives the closure of $(\sigma_1\sigma_2)^5\sigma_1^2$.

\subsection{Reproducibility}\label{ss:repro}
The programs that carry out the computations of this appendix are included as ancillary files with the arXiv
submission of this paper, and they are also available in the public repository
\begin{center}
\url{https://github.com/bwuebben/pillowcase-bounding-cochain}\,.
\end{center}
In the ancillary files they are in the directory \path{code/}, and in the repository above in its directory
\path{paper2/code/}.
Each program prints the numbers of the statements it checks next to the computed values and reports whether
the two agree, and the file \path{README.md} lists the statements that each program checks. The programs use numpy,
sympy and mpmath, and the verified computation of Remark~\ref{rem:chirality} uses SnapPy within Sage~\cite{SnapPy}.
Apart from that computation, the proofs of Theorems~\ref{thm:inat}, \ref{thm:hhk} and~\ref{thm:twisted} do not depend
on any of this code.
